% ============================================================================
% Chapter 6: Conclusion  (FYDP-II)
%
% Compiles with pdfLaTeX. Citation keys: lettucedetect, luna, ragtruth (FYDP-I
% references; replace with your keys).
% ============================================================================

\chapter{Conclusion}
\label{ch:conclusion}

\section{Summary}

Retrieval-augmented generation reduces but does not remove hallucination:
answers still contain statements that the retrieved context does not support.
The common way to detect them is to ask a large language model to judge each
answer, which is slow, costly and usually requires sending data to a cloud
service. This project set out to show that a small language model running on
ordinary hardware can do the job better and far more cheaply.

We built a hallucination detector for Bangla RAG systems. It reads the
retrieved context, the question and the answer together in a BanglaBERT encoder
with 110 million parameters, labels every word of the answer as grounded or
hallucinated, and merges flagged words into highlighted spans. It is delivered
as a Streamlit dashboard and a FastAPI backend and runs entirely offline.

On a test set of 588 answers the detector reached a word-level F1 of
$0.835 \pm 0.007$ and a span-level F1 of $0.583 \pm 0.016$ over three training
runs. It outperformed a lexical baseline (word F1 0.370), a larger multilingual
encoder, mmBERT (0.809), and two LLM judges run locally, Gemma~4 E4B (0.674)
and Gemma~4 12B (0.755). It wrongly flagged 4\% of faithful answers, against
26\% to 38\% for the LLM judges, and checked an answer in 57~ms on a CPU or
11.5~ms on a consumer GPU, about 220 times faster than the 12B judge. Training
with LoRA reduced the trained parameters 124-fold and the saved model to 3.4~MB,
at some cost in span accuracy.

The project met the requirements set in FYDP-I for model size, latency,
token-level output, span aggregation, parameter-efficient training and local
deployment. It changed two decisions from FYDP-I, for stated reasons: the
language changed from English to Bangla, which fills the gap identified in the
FYDP-I gap analysis, and the headline metric changed from answer-level to
word- and span-level F1, because a string-matching rule alone had already
exceeded the answer-level target. The whole project cost about \$16.

\section{Limitation}

\begin{itemize}
  \item \textbf{Labels not yet verified by people.} The dataset's labels were
  produced by a generator. The team is carrying out human inter-annotation, and
  the results should be re-checked against the verified labels.
  \item \textbf{Balanced test set.} Exactly half of the test answers are
  hallucinated. Real systems hallucinate less often, which would lower
  precision and raise the share of warnings that are false alarms.
  \item \textbf{Short answers.} 96\% of answers are a single sentence. Longer
  answers in which only one clause is wrong are harder and were not tested.
  \item \textbf{Weak hallucination types.} Omissions and entity replacements
  are detected least well (word F1 0.60 and 0.64); about 30\% of such
  answers are missed.
  \item \textbf{Context length.} BanglaBERT reads at most 512 tokens. The
  current contexts fit, but longer retrieved passages would be cut.
  \item \textbf{Scope of checking.} The detector checks an answer only against
  the context it is given. It cannot tell when the context itself is wrong, and
  it cannot fact-check text that has no source.
  \item \textbf{Comparison scope.} The LLM judges were local 4-bit models, each
  run once without their reasoning mode. Larger cloud models might score
  higher, at a much greater cost. The detector was not evaluated on English
  RAGTruth, so it cannot be compared directly with Luna and
  LettuceDetect~\cite{luna,lettucedetect}.
  \item \textbf{Licence.} BanglaBERT's licence allows research but not
  commercial use.
\end{itemize}

\section{Future Work}

\begin{itemize}
  \item \textbf{Human-verified evaluation.} Complete the team's inter-annotation
  of the dataset, report the agreement between annotators, and retrain and
  re-evaluate on the verified labels.
  \item \textbf{Realistic test conditions.} Build test sets with a realistic
  share of hallucinated answers, and with multi-sentence answers in which only
  one clause is wrong.
  \item \textbf{English comparison.} Train and evaluate the same pipeline on
  RAGTruth~\cite{ragtruth} with an English encoder such as ModernBERT, to
  compare directly with Luna and LettuceDetect.
  \item \textbf{Hard hallucination types.} Improve omission and entity
  replacement, for example with more training examples of these types, or with
  a hybrid in which the small model checks every answer and an LLM judge, which
  was strong on entity replacement, re-examines only the uncertain cases.
  \item \textbf{Span boundaries.} Test a conditional random field (CRF) output
  layer, which models neighbouring labels together and may draw span
  boundaries more accurately, especially for LoRA.
  \item \textbf{Longer contexts.} Handle retrieved passages beyond 512 tokens by
  splitting the context into overlapping windows, or with a long-context
  encoder pretrained on Bangla when one becomes available.
  \item \textbf{Deployment.} Merge the LoRA adapter into the base model to
  remove its small inference overhead, and integrate the API into a working
  Bangla RAG chatbot to measure the detector in real use.
  \item \textbf{Commercial use.} Train a version on a backbone with a
  permissive licence so that the system can be used commercially.
\end{itemize}
